% ===== PRÉAMBULE CANONIQUE — à coller VERBATIM en tête de CHAQUE .tex =====
\documentclass[11pt,a4paper]{article}
\usepackage[utf8]{inputenc}\usepackage[T1]{fontenc}\usepackage{lmodern}
\usepackage[french]{babel}
\usepackage{amsmath,amssymb,mathtools}\usepackage{icomma}
\usepackage[version=4]{mhchem}          % \ce{...} : équations & formules
\usepackage{siunitx}\sisetup{locale=FR} % \qty{}{}, \si{}, \ang{}
\usepackage{chemfig}                    % \chemfig{...} : structures développées
\usepackage{xcolor}\usepackage{colortbl}
\usepackage{tikz}\usepackage{pgfplots}\pgfplotsset{compat=1.18}
\usepackage[most]{tcolorbox}\usepackage{enumitem}\usepackage{array,booktabs}
\usepackage[a4paper,margin=2.2cm]{geometry}
\usetikzlibrary{arrows.meta,calc,angles,quotes,positioning,patterns,backgrounds,shapes.geometric}
\definecolor{primary}{RGB}{222,49,99}\definecolor{primarydark}{RGB}{150,22,55}
\definecolor{defcol}{RGB}{0,114,178}\definecolor{methcol}{RGB}{52,92,148}
\definecolor{excol}{RGB}{0,135,90}\definecolor{attncol}{RGB}{200,40,55}
\tcbset{boxbase/.style={enhanced,breakable,boxrule=0.9pt,arc=2.5pt,
  left=3mm,right=3mm,top=2.3mm,bottom=2.3mm,fonttitle=\bfseries,coltitle=white}}
% --- boîtes TUTORIEL (noms EXACTS, ne pas renommer) ---
\newtcolorbox{cle}[1][]{boxbase,colback=defcol!6,colframe=defcol,
  title={\textsc{À retenir}\ifx\relax#1\relax\relax\else\textmd{ : \itshape #1}\fi}}
\newtcolorbox{methode}[1][]{boxbase,colback=methcol!7,colframe=methcol,sharp corners,
  title={\textsc{Méthode}\ifx\relax#1\relax\relax\else\textmd{ --- \itshape #1}\fi}}
\newtcolorbox{exemplebox}[1][]{boxbase,colback=excol!5,colframe=excol,
  title={\textsc{Exemple}\ifx\relax#1\relax\relax\else\textmd{ : \itshape #1}\fi}}
\newtcolorbox{piege}[1][]{boxbase,colback=attncol!6,colframe=attncol,
  title={\textsc{Piège}\ifx\relax#1\relax\relax\else\textmd{ : \itshape #1}\fi}}
\newtcolorbox{remarque}[1][]{boxbase,colback=black!4,colframe=black!45,
  title={\textsc{Remarque}\ifx\relax#1\relax\relax\else\textmd{ : \itshape #1}\fi}}
% --- boîtes EXERCICE / CORRIGÉ (noms EXACTS, auto-comptées) ---
\newtcolorbox[auto counter]{exo}[1][]{boxbase,colback=primary!4,colframe=primary,
  title={\textsc{Exercice}~\thetcbcounter\ifx\relax#1\relax\relax\else\textmd{ --- \itshape #1}\fi}}
\newtcolorbox[auto counter]{corrige}[1][]{boxbase,colback=excol!4,colframe=excol,
  title={\textsc{Corrigé}~\thetcbcounter\ifx\relax#1\relax\relax\else\textmd{ --- \itshape #1}\fi}}
\newcommand{\R}{\mathbb{R}}\newcommand{\N}{\mathbb{N}}\newcommand{\Z}{\mathbb{Z}}\newcommand{\ds}{\displaystyle}
% (un \usepackage{fancyhdr}+en-tête est possible ; il est ignoré côté web)
% =========================================================================
\begin{document}
\begin{corrige}[équations ioniques : masse \emph{et} charge]
On équilibre d'abord les atomes autres que H et O, puis O par \ce{H2O}, puis H par \ce{H+} ; on
contrôle enfin la charge.
\begin{enumerate}[label=\alph*)]
  \item \ce{Cr2O7^2- + 6 Fe^2+ + 14 H+ -> 2 Cr^3+ + 6 Fe^3+ + 7 H2O}.\\
        Atomes : Cr $2=2$, Fe $6=6$, O $7=7$, H $14=14$. Charge : à gauche
        $-2+6(+2)+14(+1)=+24$ ; à droite $2(+3)+6(+3)=+24$. \checkmark
  \item \ce{MnO4^- + 5 Fe^2+ + 8 H+ -> Mn^2+ + 5 Fe^3+ + 4 H2O}.\\
        Atomes : Mn $1=1$, Fe $5=5$, O $4=4$, H $8=8$. Charge : à gauche
        $-1+5(+2)+8(+1)=+17$ ; à droite $(+2)+5(+3)=+17$. \checkmark
\end{enumerate}
\end{corrige}

\begin{corrige}[équations moléculaires à gros coefficients]
\begin{enumerate}[label=\alph*)]
  \item \ce{3 Cu + 8 HNO3 -> 3 Cu(NO3)2 + 2 NO ^ + 4 H2O}.\\
        Cu $3=3$ ; H $8=8$ ; N $8=6+2$ ; O $8\times3=24=3\times6+2+4$. \checkmark
  \item \ce{K2Cr2O7 + 14 HCl -> 2 KCl + 2 CrCl3 + 3 Cl2 ^ + 7 H2O}.\\
        K $2=2$ ; Cr $2=2$ ; O $7=7$ ; H $14=14$ ; Cl $14=2+6+6$. \checkmark
\end{enumerate}
\end{corrige}

\begin{corrige}[le kérosène de l'avion]
\begin{enumerate}[label=\alph*)]
  \item Carbone : $14\,\ce{CO2}$ ; hydrogène : $15\,\ce{H2O}$ ; oxygène à droite
        $14\times2+15=43$, soit $\tfrac{43}{2}\,\ce{O2}$. On double toute l'équation :
        \[
        \boxed{\ce{2 C14H30 + 43 O2 -> 28 CO2 + 30 H2O}}.
        \]
        Vérification — C : $28=28$ ; H : $60=60$ ; O : $86=56+30$. \checkmark
  \item Somme $=2+43+28+30=\boxed{103}$.
  \item $103\in[101;150]$.
\end{enumerate}
\end{corrige}

\begin{corrige}[diagnostic : trouver l'erreur]
\begin{enumerate}[label=\alph*)]
  \item Inventaire de la proposition $\ce{4 Zn + 10 HNO3 -> 4 Zn(NO3)2 + NH4NO3 + 2 H2O}$ :
  \begin{center}\small\renewcommand{\arraystretch}{1.2}
  \begin{tabular}{lccl}
  \toprule
  Élément & Réactifs & Produits & \\
  \midrule
  \ce{Zn} & $4$ & $4$ & \checkmark\\
  \ce{N}  & $10$ & $4\times2+2=10$ & \checkmark\\
  \ce{O}  & $30$ & $4\times6+3+2=29$ & \textcolor{attncol}{$\neq$}\\
  \ce{H}  & $10$ & $4+4=8$ & \textcolor{attncol}{$\neq$}\\
  \bottomrule
  \end{tabular}
  \end{center}
  L'\textbf{hydrogène} ($10\to8$) et l'\textbf{oxygène} ($30\to29$) ne sont pas équilibrés.
  \item Les deux manques (2 H et 1 O) se comblent d'un coup avec \textbf{une} molécule d'eau
        supplémentaire : le coefficient de \ce{H2O} doit être $3$, non $2$ :
        \[
        \boxed{\ce{4 Zn + 10 HNO3 -> 4 Zn(NO3)2 + NH4NO3 + 3 H2O}}.
        \]
        Hydrogène : $10=4+6$. \checkmark
  \item Oxygène : à droite $4\times6+3+3=30$, égal aux $30$ des réactifs (\ce{10 HNO3}). \checkmark
\end{enumerate}
\end{corrige}

\end{document}
